\documentclass[12pt, a4paper]{article}

% --- Pacchetti essenziali ---
\usepackage[utf8]{inputenc}
\usepackage[T1]{fontenc}
\usepackage[italian]{babel}
\usepackage{geometry}
\geometry{margin=2.5cm}

% --- Grafica e figure ---
\usepackage{graphicx}
\usepackage{float}

% --- Matematica ---
\usepackage{amsmath}
\usepackage{amssymb}

% --- Codice sorgente ---
\usepackage{listings}
\usepackage{xcolor}
\lstset{
    language=C++,
    basicstyle=\ttfamily\small,
    keywordstyle=\color{blue}\bfseries,
    commentstyle=\color{gray}\itshape,
    stringstyle=\color{red},
    numbers=left,
    numberstyle=\tiny\color{gray},
    breaklines=true,
    frame=single,
    captionpos=b
}

% --- Link e riferimenti ---
\usepackage{hyperref}
\hypersetup{
    colorlinks=true,
    linkcolor=black,
    citecolor=black,
    urlcolor=blue
}

% --- Bibliografia ---
\usepackage{biblatex}
\addbibresource{bibliography.bib}

% -------------------------------------------------------
% FRONTESPIZIO
% -------------------------------------------------------
\title{
    \vspace{2cm}
    {\Large Università di Bologna} \\
    \vspace{0.5cm}
    {\large Corso di ...} \\
    \vspace{2cm}
    \rule{\linewidth}{0.5pt} \\
    \vspace{0.5cm}
    {\LARGE \textbf{Titolo della Relazione}} \\
    \vspace{0.3cm}
    {\large Sottotitolo (opzionale)} \\
    \vspace{0.5cm}
    \rule{\linewidth}{0.5pt}
}
\author{
    Nome Cognome \\
    \small Matricola: XXXXXX \\
    \small \texttt{email@studenti.unibo.it}
}
\date{Anno Accademico 2025/2026}

% -------------------------------------------------------
\begin{document}

\maketitle
\thispagestyle{empty}
\newpage

% --- Indice ---
\tableofcontents
\newpage

% -------------------------------------------------------
% ABSTRACT
% -------------------------------------------------------
\begin{abstract}
Breve descrizione del progetto o dell'elaborato (5-10 righe). 
Riassumi l'obiettivo, il metodo utilizzato e i risultati principali.
\end{abstract}

\newpage

% -------------------------------------------------------
% 1. INTRODUZIONE
% -------------------------------------------------------
\section{Introduzione}
Descrivi il contesto del problema affrontato. Spiega:
\begin{itemize}
    \item Qual è il problema o l'obiettivo
    \item Perché è rilevante
    \item Come è strutturata la relazione
\end{itemize}

% -------------------------------------------------------
% 2. BACKGROUND E TEORIA
% -------------------------------------------------------
\section{Background e Teoria}
Esponi i concetti teorici necessari per comprendere il lavoro.
Puoi inserire formule matematiche:
\begin{equation}
    E = mc^2
\end{equation}

% -------------------------------------------------------
% 3. METODOLOGIA
% -------------------------------------------------------
\section{Metodologia}
Descrivi come hai affrontato il problema:
\begin{itemize}
    \item Algoritmi o tecniche utilizzate
    \item Strumenti e librerie usate
    \item Scelte progettuali
\end{itemize}

% -------------------------------------------------------
% 4. IMPLEMENTAZIONE
% -------------------------------------------------------
\section{Implementazione}
Descrivi i dettagli implementativi. Puoi inserire snippet di codice:

\begin{lstlisting}[caption={Esempio di codice C++}]
#include <iostream>

int main() {
    std::cout << "Hello, World!" << std::endl;
    return 0;
}
\end{lstlisting}

% -------------------------------------------------------
% 5. RISULTATI
% -------------------------------------------------------
\section{Risultati}
Presenta i risultati ottenuti. Puoi inserire immagini:

% \begin{figure}[H]
%     \centering
%     \includegraphics[width=0.7\textwidth]{immagine.png}
%     \caption{Didascalia immagine}
%     \label{fig:etichetta}
% \end{figure}

% -------------------------------------------------------
% 6. DISCUSSIONE
% -------------------------------------------------------
\section{Discussione}
Analizza i risultati:
\begin{itemize}
    \item Cosa funziona bene
    \item Limitazioni riscontrate
    \item Possibili miglioramenti
\end{itemize}

% -------------------------------------------------------
% 7. CONCLUSIONI
% -------------------------------------------------------
\section{Conclusioni}
Riassumi il lavoro svolto e i risultati principali. 
Indica possibili sviluppi futuri.

% -------------------------------------------------------
% BIBLIOGRAFIA
% -------------------------------------------------------
\newpage
\printbibliography[title={Bibliografia}]

\end{document}